\section{后汉后周的军队、税源与淮南接收}
\label{sec:later-five-dynasties-armies-tax}

\noindent	extbf{战争不是地图上的箭头。} 948--960年的后汉继承危机、后周与北汉、高平、寺院整顿、南征淮南、江北接收、北伐与陈桥，都要求把军事行动放回粮道、马匹、舟船、城防、临时征发和信息传递中理解。军队能够在某月出发，不等于后方已有稳定供给；一座城被攻取，也不等于城内人口、市场和文书立刻服从新主人。正史传记常将胜负归于将领的决断，而实际的行军还取决于河流是否可渡、仓粮能否转运、地方军是否愿意协同、居民又以何种方式躲避或应对。这些条件很少被保存为完整账簿，正需要叙事避免把军事数字和报功语句变成未经检验的事实。

\noindent	extbf{财政的现场。} 战争时期的税、盐、漕粮、寺院资产、市场关津和地方留截，都是维持武装集团的现实条件。朝廷试图通过诏令、使职、户籍与任命把资源汇集到中枢，军镇和区域政权则会以同样的文书语言巩固自己的军赏、人事和仓储。研究者能从制度志确认政策的名目，从地方文书看见个别纳税、借贷或役使，却不能将任一材料直接扩大成全国收支表。更稳妥的判断是追踪资源流向的改变：哪些道路被反复争夺，哪些城镇成为粮运节点，哪些收入被地方军事化，哪些家庭因战争而退出原有名册。

\noindent	extbf{区域社会的主动性。} 地方并非只在中央崩解后才拥有历史。家族、寺院、商人、书吏、工匠、军户和流动者早已依靠亲属、市场、信仰与地理维持生活；政权转换改变的是这些关系被征用、保护和书写的方式。墓志和碑刻使精英的迁徙与任官特别可见，契约和愿文能偶见较低层次的财产和宗教活动，遗址则显示生产或交通的物质条件。它们不能替沉默的大多数发言，却足以拒绝将地方社会化约为接受诏令或等待征服的被动背景。

\noindent	extbf{边地与多重时间。} 草原、河西、岭南、淮南、巴蜀和海港各有不同的季节、路程与政治语言。册封、和盟、置州、投降或受禅，是可以由材料确证的动作，不是统治深度的同义词。汉文朝廷记录强调秩序与朝贡，邻国史书、碑铭和多语文书又有不同的观察位置；材料不一致时，应保留这种不一致，而不是强行归入单一国家叙事。多政权并立时期尤其如此：并行年号和相互册命既是竞争，也是维持贸易、使节和地方人事的工具。

\noindent	extbf{年度细读。} 账本编号LT-948-108、LT-951-109、LT-951-110、LT-954-111、LT-955-112、LT-956-113、LT-958-115、LT-960-118为本节提供逐年可复查的入口。年度锚点并不自动含有一整年的社会史，却能防止后来的结局吞没当时的选择。把同一年内的中枢诏令、跨政权行动和制度社会材料并置，才能看见某些地区在战争中恢复、另一些地区在和平语汇下继续被征发的不同节奏。对755--960年的解释，因而应当是多条时间线的交织，而不是一条由盛而衰、再由乱而治的单向曲线。

\noindent	extbf{证据与限度。} 本节主要依托《旧五代史》《五代会要》《辽史》《宋史》《资治通鉴》及江淮遗址。官修史书、制度汇编、文书、碑刻和考古遗存各自照亮不同问题：前者擅长年序和中央语言，后者更接近局部实践，却也有保存偶然性和自我表述。可以确认政策发布、军队移动、城市易手和名号转换；不能由此断言所有地区同步执行、所有人共享同一动机，或把一条材料中的数字视为精确统计。把限度写入叙事，不是削弱判断，而是让军政、财政、区域社会和文化记忆在可承受的证据范围内相互关联。

\subsection{军队、家庭与地方秩序}
军队的存在总会重新排列地方社会。驻军需要粮、草、房舍、修具和向导，征调会把青壮年、牲畜和船只从原有生产中抽走；但不同人群并不以同一种方式回应。拥有土地、店铺或书写能力的家庭可能试图借官府凭据保护财产，流动者和贫困者则更可能被征入军中、依附强者、迁往他处，或从现存记录中消失。战争史若只追随主将与都城，便会遗漏这种不平等的承受。墓志中一次迁葬、一份契约中的担保人、寺院愿文里的施主名单，都不足以恢复全体人口，却能提示人们，所谓地方安定是由许多不均等的安排暂时拼合出来的。

军镇与区域政权的内部也并非铁板一块。将领要依靠亲兵、部将、文吏和地方豪强，才能把盐利、田赋或关津收入转化为军赏；部众又会在继承、欠饷和作战失利时重新选择归属。因此，“节度使据地”或“某国称帝”只是解释的起点，不是对军政结构的终点。应继续追问名义任命是否进入州县，人事变动是否改变仓储和道路，军队是否拥有稳定补给，中央或邻国的使节关系又为谁创造了议价空间。这些问题能把抽象的割据、统一和边防还原为可观察的组织劳动。

\subsection{交通、市场与物质条件}
道路和水路不是被动承受战争的背景。河流涨落、山道险易、港口停泊条件、驿站修缮和城池仓廪，直接决定军队可否迅速移动，也决定商人、移民、僧侣和使节能否跨区域往来。交通改善并不自动带来共同繁荣：同一条水路可以运送粮食救济，也可以将地方余粮优先输往军营；同一座港市可以连接远方市场，也可能在战乱中成为掠夺和征税的目标。考古遗存能说明某条线路或某类生产曾经存在，尚不足以说明每次运输的规模和受益者。将地理限制写进叙事，正是为了避免把国家决策描写成脱离物质世界的意志。

货币、绢帛、盐、茶、金属和寺院财产在这一时期不断进入军政计算，却不宜被理解为彼此独立的“经济部门”。铸钱需要铜料和劳力，军赏需要可被接受的支付形式，盐税和商税依赖市场与交通，寺院资产则牵连施主、僧侣、土地和地方声望。国家的财政命令可以试图重新分类这些资源，地方社会的交换和规避却使执行始终带有区域差异。材料保存得越零散，越应拒绝用单一数字描绘所有区域的富裕或贫困。

\subsection{权力的文书形式}
诏书、榜文、盟约、册命、契约、账簿和碑刻构成权力得以被传递和被争论的媒介。文书不仅记录已经完成的统治，也参与制造统治：它为征收定名，为军镇授予合法语言，为投降和受禅安排程序，为家族留下可诉诸的记忆。可是，文书的存在从不保证命令已被执行。抄写、传递、解释和保存都可能改变其效果；许多口头协商、临时强制和地方抵抗则根本没有进入纸面。史家应当同时把文书当作证据和行动，而不是透明的事实窗口。

\subsection{宗教与战争记忆}
战争并未使宗教生活停止。寺院可成为施舍、避难、抄写、储藏和地方结社的空间，也会因土地、劳力、金属和政治声望而进入财政与军政竞争。官方的保护、限制或废并只能说明国家的目标，不能概括所有寺院和信众的经历；僧传强调高僧与神异，碑刻强调功德与施主，二者都需要同地方经济和交通条件一起阅读。宗教材料尤其提醒我们，战乱中的社会并非只有破坏，也有记忆、救济、重建和重新组织关系的实践。

\subsection{避免目的论}
从唐亡看到宋兴，或从后晋的燕云安排看到后来的边防困境，都会诱使叙事把后来结果写成当时唯一可能的道路。年度材料提供的抵抗正在于：每一次兵变、降附、册封、联盟和征税，都发生在行动者尚不知结局的条件下。把不确定性留在文本中，既不否认结构限制，也不把地方社会、边地政权和普通人的选择删成单一结局的注脚。

\subsection{比较的尺度}
同一项政策在不同区域会产生不同后果。靠近都城的州县、掌握水运的江淮、拥有边军的河北、连接绿洲的河西、面向海贸的岭南或两浙，都有不同的土地、道路与人事条件。因而不能把中枢的一个诏令直接读成全国社会的时间表，也不能用某一区域的文书代替所有地区。比较并不要求材料完全相同；它要求写作者说明材料来自哪里、保存了谁、遗漏了谁，以及为什么这一差异会改变解释。

地方精英的作用也须放在这种尺度里理解。他们可能担任军府文吏、参与漕运、出资修寺、经营市场或以墓志维护家族名望；这些行动使他们成为国家与地方之间的中介，却不代表他们能够代表全体居民。中介地位既能带来资源，也会在战乱、继承和财政重组中造成风险。对精英材料的认真阅读，应当同时包含对其局限的说明。

\subsection{环境与风险}
气候、疫病、河流和地形并不决定政治结局，却会提高或降低行动的成本。远征遇到疾病和风浪、城防依赖水源、漕运受河道影响、边军要面对草场与冬季，这些条件使同一份军事命令在不同季节有不同结果。史书中的灾异语言带有政治修辞，不能直接说明因果；但完全忽略物质环境，也会让战争和财政变成脱离身体、土地与运输的抽象故事。

\subsection{持续的问题}
本节所展示的并非完整社会统计，而是一组可继续检验的问题：谁在何时能够获得粮食和安全，谁被纳入征发，谁借助文书和网络保存财产，谁在材料里没有声音。随着新的文书、考古与区域研究加入，具体答案可以修正；对材料偏向和区域差异的警觉应当保留。


\noindent 叙事的责任正在于把这样的差异写清，而不以方便的概括抹去道路、劳力、性别、身份和地区造成的不均等。可被确证的事件应当同可被合理推测的条件分开，缺少证据的地方则明确保留为问题。

材料并不要求我们放弃判断，却要求每一个判断说明它依靠何种记录、覆盖何处、又遗漏何人。只有如此，军事财政和区域社会才能同时成为历史解释的对象。

证据的沉默也是这一段历史必须面对的事实。

历史解释必须同材料的边界并存。

仍须持续核查。
